\par\noindent\SourceLeaf{52}%
\hspace*{\parindent}B) Suchen wir nun allgemein das Gesetz für den Koeffizienten, den der Kosinus eines ungeraden Vielfachen, zum Beispiel $\cos(2q-1)\varphi$, erhält. Unter den Gliedern der Reihe $(1+n\cos\varphi)^{\nu}$, die nach Potenzen von $\cos\varphi$ fortschreitet, enthält das erste mit $\cos\varphi^{2q-1}$ bei der Entwicklung nach Kosinus von Vielfachen das gesuchte Vielfache in seinem $(q-1)^{\text{to}}$ Glied. Nach dem zuvor in $(2,\beta)$ dargelegten Gesetz hat nämlich die Entwicklung von $\cos\varphi^{2q-1}$ als $(q-1)^{\text{tus}}$ Glied

\[
= \frac{1}{2^{2q-2}} {}^{2q-1}\mathfrak{B}^{(2q-1)} \cos(2q-1)\varphi,
\]

\noindent oder, da ${}^{2q-1}\mathfrak{B}^{(2q-1)} = 1$ gilt,

\[
\frac{1}{2^{2q-2}} \cos(2q-1)\varphi..
\]

\noindent In der Grundreihe war aber das Glied, das $\cos\varphi^{2q-1}$ enthält, mit ${}^{\nu}\mathfrak{B}^{(2q-1)} n^{2q-1}$ multipliziert. Fügt man diesen Faktor hinzu, so erhält man den ersten Teil des zu $\cos(2q-1)\varphi$ gehörenden Koeffizienten oder das Anfangsglied der Reihe, die diesen Koeffizienten angibt

\[
= \frac{1}{2^{2q-2}} {}^{\nu}\mathfrak{B}^{(2q-1)} n^{2q-1}.
\]

\noindent Jede ungerade Potenz von $\cos\varphi$, deren Exponent größer als $2q-1$ ist, liefert nun, wenn man sie durch Kosinus von Vielfachen ausdrückt, ein mit $\cos(2q-1)$ multipliziertes Glied. Das $h^{\text{tus}}$ unter diesen entsteht aus der $h^{\text{ta}}$ ungeraden Potenz nach $\cos\varphi^{2q-1}$, deren Exponent also $2q+2h-1$ ist. Man nimmt aus der Größe

\[
{}^{\nu}\mathfrak{B}^{(2q+2h-1)} n^{2q+2h-1} \cos\varphi^{2q+2h-1}
\]

\noindent nach Entwicklung in Kosinus von Vielfachen das $(q-1)^{\text{tus}}$ Glied. Dieses ergibt sich wie folgt (da nach dem bekannten Gesetz für $\cos\varphi^{2q+2h-1}$ das $(q-1)^{\text{tus}}$ Glied lautet

\[
= \frac{1}{2^{2q+2h-2}} {}^{2q+2h-1}\mathfrak{B}^{(2q+h-1)} \cos(2q-1)) 
\]

\[
= \frac{1}{2^{2q+2h-2}} {}^{2q+2h-1}\mathfrak{B}^{(2q+h-1)}  .  {}^{\nu}\mathfrak{B}^{(2q+2h-1)} n^{2q+2h-1} \cos(2q-1)\varphi.
\]

\noindent Der Koeffizient dieser Größe ist also der $h^{\text{tus}}$ unter denen, die $\cos(2q-1)\varphi$ bei der aufeinanderfolgenden Entwicklung der Glieder der Grundreihe erhält. Fassen wir daher bei der Bildung der neuen Reihe die Summe dieser Teile, wie es zulässig ist, als einen einzigen Koeffizienten auf, so ist der Koeffizient von $\cos(2q-1)\varphi$ gleich einer Reihe mit dem Anfangsglied

\[
= \frac{1}{2^{2q-2}} {}^{\nu}\mathfrak{B}^{(2q-1)} n^{2q-1}
\]
